% !TeX program = xelatex
\documentclass{kkubee_edit_all}

\usepackage{amsthm}
\usepackage{amsfonts}
\usepackage{amssymb}
\usepackage{mathtools}

%% ============ Theorem Environments ============
\newtheorem{thrm}{ทฤษฎีบท}[chapter]
\newtheorem{ex}[thrm]{ตัวอย่างที่}
\newtheorem{rem}[thrm]{หมายเหตุ}
\newtheorem{warning}[thrm]{คำเตือน}
\newtheorem{prop}[thrm]{บทแทรก}
\newtheorem{cor}[thrm]{บทสรุป}
\newtheorem{definition}[thrm]{นิยาม}
\newtheorem{exercise}{แบบฝึกหัด}[chapter]

%% โหลด mathstyle ที่เราสร้างไว้
\usepackage{mathstyle}

\begin{document}
\chapter*{สาธิตรูปแบบกล่องคณิตศาสตร์ (Modern Math Style Demo)}
\addcontentsline{toc}{chapter}{สาธิตรูปแบบกล่องคณิตศาสตร์}

เอกสารนี้แสดงตัวอย่างรูปแบบกล่องข้อความคณิตศาสตร์ต่างๆ ที่กำหนดไว้ใน \texttt{mathstyle.sty} ซึ่งได้รับการปรับปรุงให้มีสีสันโทน Pastel ดูทันสมัย ขอบมน และมีเงา (Drop shadow) เพื่อให้ผู้เรียนอ่านง่ายและสบายตา

\vspace{1cm}

\begin{definition}[นิพจน์ (Expression)]
ในทางคณิตศาสตร์ นิพจน์คือการรวมกันของตัวเลข ตัวแปร และตัวดำเนินการ ที่สามารถหาค่าได้ ตัวอย่างเช่น $3x + 5$ หรือ $y^2 - 2y + 1$
\end{definition}

\begin{thrm}[ทฤษฎีบทพีทาโกรัส]
สำหรับรูปสามเหลี่ยมมุมฉากใดๆ พื้นที่ของรูปสี่เหลี่ยมจัตุรัสบนด้านตรงข้ามมุมฉาก เท่ากับผลรวมของพื้นที่ของรูปสี่เหลี่ยมจัตุรัสบนด้านประกอบมุมฉากทั้งสอง
\[ c^2 = a^2 + b^2 \]
\end{thrm}

\begin{prop}[สมบัติการสลับที่]
การบวกและการคูณมีสมบัติการสลับที่ ดังนี้
\begin{itemize}
    \item $a + b = b + a$
    \item $a \cdot b = b \cdot a$
\end{itemize}
\end{prop}

\begin{cor}
จากทฤษฎีบทก่อนหน้า สามารถสรุปได้ว่าสมการเชิงเส้นตัวแปรเดียวมีเพียงคำตอบเดียวเสมอ
\end{cor}

\begin{rem}
\text{หมายเหตุ:} ตัวแปรที่ใช้ในสมการนี้ต้องเป็นจำนวนจริง (\text{Real Numbers}) เท่านั้น
\end{rem}

\begin{warning}[ข้อควรระวัง]
ห้ามหารด้วยศูนย์ ($0$) ในทุกกรณี เนื่องจากไม่มีนิยามในทางคณิตศาสตร์ และจะทำให้เกิดข้อผิดพลาดในการคำนวณ
\end{warning}

\begin{ex}[การแก้สมการเชิงเส้น]
จงหาค่าของ $x$ จากสมการ $2x + 5 = 15$
\end{ex}

\begin{solution}
\begin{solsteps}
    \item จากสมการ $2x + 5 = 15$
    \item นำ $5$ ไปลบทั้งสองข้าง: $2x = 15 - 5$
    \item จะได้ $2x = 10$
    \item นำ $2$ ไปหารทั้งสองข้าง: $x = \frac{10}{2}$
    \item ดังนั้นคำตอบคือ \mathoval{x = 5}
\end{solsteps}
\end{solution}

\begin{exercise}
จงพิสูจน์ว่าหาก $x^2$ เป็นจำนวนคู่ แล้ว $x$ จะต้องเป็นจำนวนคู่
\end{exercise}

\begin{proof}
สมมติให้ $x$ เป็นจำนวนคู่ ดังนั้น $x = 2k$ สำหรับจำนวนเต็ม $k$ บางตัว
เมื่อนำมาพิจารณา $x^2$:
\[ x^2 = (2k)^2 = 4k^2 = 2(2k^2) \]
เนื่องจาก $2k^2$ เป็นจำนวนเต็ม จะได้ว่า $x^2$ เป็นจำนวนคู่ \qedhere
\end{proof}

\end{document}
